\subsection{从山海关到北京：一座城如何变成朝廷}

顺治元年四月，山海关把三种彼此利用、也彼此不信任的力量挤到了一起。李自成进入北京后，关宁军同大顺在兵饷和统属上已经决裂；多尔衮则把吴三桂的处境变成八旗越关的机会。会战打开的是通往华北的道路，不是天下已经归属的证明。五月入北京之后，清军首先要接收的是仓储、官署、粮道和等待观望的官员；冬季的招纳、征发与整顿官署，才试图把一支行军的军队变成能收发文书的政权。

这份接收有很具体的物质尺度。驻军的口粮、马匹的草料、南下所需的船只和车夫，都不能由北京城门已经易帜而自动出现；而搜集这些资源又会把新政权带到州县的仓廒、渡口和里甲面前。对朝廷而言，旧有的赋役名目、漕运节点和能办案的书吏，是让军令继续向外走的条件；对被征发的人和被留任的官吏而言，它们首先意味着饷额、差役与人身安全尚未有定数。故而“接收明制”既是制度上的借用，也是一次围绕粮、船、人力与地方信誉不断谈判的过程。

北京并非唯一的政治空间。顺治帝迁驻北京，使皇帝、部院和礼制中心汇合；盛京仍保有旗务、皇室与王朝记忆的重量。与此同时，江南的水路、两广和西南的山道、福建沿海的港口与岛屿，各有不同的通行条件、供给办法和对手。摄政政权能在京师发出同一格式的命令，却不能使这些区域以同一速度进入其财政和军政网络。把清初写成由北京向四方均匀铺开的统治，会把多区域的战争误写成一张已经完成的行政地图。

摄政结束、顺治亲政，也没有消除这种多重结构。北京需要熟悉赋役、漕运与刑名的书吏和降官，前线持续索要粮饷、船只和兵丁；皇帝的诏令必须经过宗室、旗务、满汉官僚、将领以及地方中介，才可能抵达一州一县。康熙幼年时的四辅政大臣、鳌拜被拘捕和皇帝重新掌握人事，显示宫廷权力的重新分配；它们并不证明命令从宫门笔直地变成地方的服从。每一次任命与调度仍受道路、军费、旧属关系和战区局势牵制。

征服的代价也写在城池和身体上。南京失守以后，弘光、隆武、绍武、永历等政治中心仍在东南和西南移动；一座府城易手不等于乡村、岛屿、山道和税源同时易手。薙发易服的命令把政治归属做成可见的标记，却在各地同守城、征发、逃亡和报复纠缠。有人借服从求得暂时的秩序，有人因旧主、亲属、粮饷或城防而延宕、离散或抵抗；这些选择不能被压成一条从命到不从命的整齐线。

扬州、江阴和广州的杀戮与损毁不能被胜利叙事抹去，但也不宜以一个貌似精确的总数代替不同地方的经历。城守记述、幸存者叙事和后修方志使围城、屠戮、逃亡与毁坏可见；它们又常带有传抄、追忆、纪念共同体和数字修辞的痕迹。清廷战报所记的攻取与招降，也不是伤亡普查。能够确定的是征服暴力深刻改变了城市和周边社会；不能由现存文本合算一项代表全境的死亡数字。直到永历政权覆亡、夔东据点被攻破，郑氏仍据海上与台湾，所谓统一始终是多战区逐段取得的结果。

\subsection{材料与争议}

《清世祖实录》、内国史院满文档案和题本能追踪北京政权怎样制作命令、接收官署和调兵筹饷；它们首先是新政权的文书，而非命令在远方已经实现的证明。满、汉文本也不能未经核对便互作替身：官号、地名和分类用语的对应，会影响我们所见的行动者与行政范围。题本所留下的是呈报、批示和案卷得以进入中枢的链条，尤其难以覆盖未能上报的逃散、拖延和私下交涉。

南明奏疏、地方记事、城守记录与后出的方志保存败亡朝廷和城市共同体的不同时间，也有传抄、追忆和地方立场。它们与官书相对读，不是为了拼出一种全知的中性叙事，而是为了分别检验城池何时易手、何种征发可见、哪些人被写入或留在沉默中。山海关的求援文书、各军兵数与攻城伤亡因此不能被写成无争议的精确场景；同样，中央命令不能代替地方接受度的统计。能够确定的是征服的暴力、反复与制度接收同时发生；各地的服从、损失与恢复程度则须随材料的保存状况分别判断。

\subsection{交叉阅读窗口：一纸薙发谕如何把征服写进身体}

顺治二年六月谕旨以“遵依者为我国之民，迟疑者同逆命之寇”划分服从者与敌人。短句出自《清世祖实录》的朝廷命令，和内国史院档案一样，首先可证明新政权如何发出、归类和威吓；它的对偶和二分法把发式压缩成政治归属，因而具有比行政术语更强的逼迫力。它不是江南每一处执行情形的逐日记录，更不能说明每个遵从者都出于同一种判断。命令越过城门、渡口和驿路时，会同驻军补给、城防崩解、地方差役与家庭避难交织在一起。

钱穆会追问的不是这句谕令是否“有效”，而是摄政政权怎样借明代官署、旗务与礼制把命令变成可持续的政治关系；吕思勉式的问题则把视线移向军饷、赋役、家庭逃散与城镇社会如何限制或改变这种关系。两种提问都不能替代证据：实录和档案证明命令与中央分类，地方记录仅能局部检验收到、抵抗或服从的过程。这个窗口因此不把一句严厉的国家声音误作征服已经完成的证明。
